\documentclass[a4paper, 11pt,twoside]{article}
\usepackage[utf8]{inputenc}
\usepackage[T1]{fontenc}
\usepackage{lmodern}
\usepackage[french]{babel}
\usepackage{amsmath, amssymb, amsthm}
\usepackage{geometry}
\usepackage{graphicx}
\usepackage{enumitem}
\usepackage{hyperref}
\usepackage{siunitx}
\usepackage{booktabs}
\usepackage{fontawesome} % Pour des icônes modernes comme des étoiles
\usepackage{multirow}
\usepackage{float}
\usepackage{fancyhdr}
\usepackage{tcolorbox} % Pour les encadrés avec bords arrondis
%\usepackage{fontspec} % Pour utiliser des polices personnalisées
\usepackage{tikz} % Pour des éléments graphiques
\usepackage[tikz]{bclogo}
\usepackage{titlesec} % Pour personnaliser les sections et sous-sections
\usepackage{fmtcount} % Pour formater les numéros en lettres
\usepackage{qrcode}


% Commande pour lien vers le corrigé
\newcommand{\lcorr}[1]{~\hfill
\textit{Corrigé : }~
\begin{minipage}{10em}
\qrcode[height=10em]{#1}
\end{minipage}}

% Commande pour lien absolu vers les images
\newcommand{\includegraphic}[2]{
\includegraphics[width=#1\linewidth]{C:/Users/simon/Google Drive/Enseignement/CPGE/MPE Wallon/Latex/Exercices/#2}
}




% Configuration de la géométrie de la page
\geometry{hmargin=2.5cm, vmargin=2.5cm}


% Charger le package xcolor pour les couleurs
\usepackage{xcolor}
\definecolor{titleblue}{RGB}{0, 0, 255} % Définir une couleur bleue
\definecolor{vertclair}{RGB}{144, 238, 144} % Vert clair
\definecolor{bleu}{RGB}{0, 0, 255} % Bleu
\definecolor{rouge}{RGB}{255, 0, 0} % Rouge
\definecolor{noir}{RGB}{0,0,0}% Noir
\definecolor{grisclair}{rgb}{0.9,0.9,0.9}% Noir

% Supprimer l'alinéa en début de paragraphe
\usepackage{parskip}
\setlength{\parindent}{0pt} % Pas d'alinéa
\setlength{\parskip}{0.5em} % Espacement entre les paragraphes

% Style des titres avec encadré et bords arrondis
\usepackage{titlesec}


% Définir une commande pour créer un tcolorbox avec le titre
\newcommand{\titlebox}[1]{%
  \begin{tcolorbox}[colback=gray, colframe=black, arc=0mm, boxrule=0pt]
     \textcolor{white}{\textbf{#1}} % Titre en gras

  \end{tcolorbox}%
}
\newcommand{\subsubsectionbox}[1]{%
  \begin{tcolorbox}[colback=grisclair, colframe=black, arc=0mm, boxrule=0pt]
     \textcolor{black}{\textbf{#1}}
  \end{tcolorbox}%
}

% Redéfinir la numérotation des sections en lettres (A, B, C, ...)
%\renewcommand{\thesection}{\Alph{section} -}
%
% Redéfinir la numérotation des sous-sections en "Exercice 1 :", "Exercice 2 :", ...
%\renewcommand{\thesubsubsection}{Exercice \arabic{subsubsection} :}




% Commande personnalisée pour les titres de sous-sections avec tcolorbox
\newcommand{\subsubsectiontitle}[1]{%
  \subsubsectionbox{\thesubsubsection\ #1}%
}


\titleformat{\section}
  {\normalfont\Large\bfseries\sffamily} % CMU Sans Serif en gras
  {\thesection}{}{\titlebox}

\titleformat{\subsection}
  {\normalfont\large\bfseries\sffamily} % CMU Sans Serif en gras
  {\thesubsection}{1em}{}

\titleformat{\subsubsection}
  {\normalfont\normalsize\bfseries\sffamily} % CMU Sans Serif en gras
  {\thesubsubsection}{1em}{\subsubsectionbox}



%\titleformat{\subsubsection}
%  {\normalfont\normalsize\bfseries\sffamily} % CMU Sans Serif en gras
%  {\thesubsubsection}{1em}{} % Pas de boîte pour les sous-sous-sections
%

% Commandes pour les carrés de couleur dans la marge gauche
\usepackage{tikz} % Pour dessiner les carrés
\usepackage{marginnote} % Pour placer les carrés dans la marge

% Définir les carrés de couleur
\newcommand{\facile}{%
  \marginnote{%
    \begin{tikzpicture}[baseline=-0.5ex]
      \fill[vertclair] (0,0) rectangle (0.5,0.5); % Carré vert clair
    \end{tikzpicture}%
  }%
}
\newcommand{\moyen}{%
  \marginnote{%
    \begin{tikzpicture}[baseline=-0.5ex]
      \fill[bleu] (0,0) rectangle (0.5,0.5); % Carré bleu
    \end{tikzpicture}%
  }%
}

\newcommand{\ouvert}{%
  \marginnote{%
    \begin{tikzpicture}[baseline=-0.5ex]
      \fill[noir] (0,0) rectangle (0.1,0.5); % Carré bleu
    \end{tikzpicture}%
  }%
}
\newcommand{\dur}{%
  \marginnote{%
    \begin{tikzpicture}[baseline=-0.5ex]
      \fill[rouge] (0,0) rectangle (0.5,0.5); % Carré rouge
    \end{tikzpicture}%
  }%
}

% Ajuster la marge gauche pour les carrés
\reversemarginpar % Placer les notes dans la marge gauche


% Ajuster l'espacement sous les sections
\titlespacing{\section}{0pt}{3em}{-1em} % Espacement avant et après la section
\titlespacing{\subsubsection}{0pt}{2em}{-1em} % Espacement avant et après la sous-section
%\titlespacing{\subsubsection}{0pt}{1em}{0.5em} % Espacement avant et après la sous-sous-section

% Style des listes
\usepackage{enumitem}
\setlist[enumerate,1]{label=\textbf{Q\arabic*.}, % Format des numéros de questions
                     leftmargin=3.2em, % Alignement à gauche
                     itemsep=0em,  
		  align=left, % Alignement à gauche
                     labelwidth=2.8em, % Largeur fixe pour les labels
                     labelsep=0.5em, % Espace entre le label et le texte% espace entre les items
                     topsep=0pt} % Pas d'espace avant la liste

% Style des équations
\newcommand{\eq}[1]{\begin{equation}#1\end{equation}}
\newcommand{\eqn}[1]{\begin{equation*}#1\end{equation*}}

% Style des tableaux
\renewcommand{\arraystretch}{1.2}

% Style des figures
\newcommand{\fig}[3]{
  \begin{figure}[H]
    \centering
    \includegraphics[width=#2]{#1}
    \caption{#3}
  \end{figure}
}

% Difficulté des Exercices
\newcommand{\difficulte}[1]{%
  \hfill\textcolor{black}{\textbf{\large%
    \ifnum#1=1\faStar\fi%
    \ifnum#1=2  \faStar\hspace{-0.1em}\faStar\fi%
    \ifnum#1=3 \faStar\hspace{-0.1em}\faStar\hspace{-0.1em}\faStar\fi%
  }}%
}



% Logo pour corrigé ou indications
\newcommand{\corrige}{ \textcolor{blue}{\faFileTextO}} % Logo pour un corrigé disponible
\newcommand{\indications}{ \textcolor{orange}{\faLightbulbO}} % Logo pour des indications disponibles
\newcommand{\lesdeux}{ \textcolor{orange}{\faLightbulbO}\hspace{-0.1em}  \textcolor{blue}{\faFileTextO}} % Logo pour les deux disponibles

\newcommand\gauchedroite[2]{\noindent{}#1\hfill#2}


% Texte pour thème et provenance
\newcommand{\sousinfo}[1]{%
  \vspace{-1.2em} % Espacement après la sous-section
  {\itshape\footnotesize #1} % Texte en italique et petit
  \vspace{1.2em} % Espacement avant le texte suivant
}

%\newenvironment{gbar}[1]{ %
%\def \FrameCommand{{\color{#1}\vrule width=3pt}\colorbox{fbase}} %
%\MakeFramed{\advance \hsize-\width   \FrameRestore }} %
%{\endMakeFramed}

\renewcommand{\headrulewidth}{0pt}

\newcommand{\planche}[2]{
\begin{tcolorbox}[colback=white,colframe=black!75!black]
\gauchedroite{Physique-Chimie }{#1 Wallon}\par
\begin{center}
    \LARGE{\textbf{Interrogation Orale}}

\vspace{0.5em}

\large{\textbf{Planche #2}}
\vspace{0.5em}
\end{center}
\end{tcolorbox}
}